\section{A quantitative library for analysis and approximation}
\label{sec:analysis-library}

This section records the shared analytic lemmas that should be formalized once
and reused.  The purpose is not to recreate a set-theoretic real-analysis
library inside iPRA.  It is to isolate finite rational contracts sufficient for
the theorem packets.

\subsection{Dyadic and rational names}

A fast rational name is a fixed evaluator \(a(p,k)\in\Q\) satisfying
\begin{equation}
 |a(p,k)-a(p,\ell)|\le 2^{-k}+2^{-\ell}.
 \label{eq:regular-name}
\end{equation}
Rationals are coded by signed numerators and positive denominators, or by
normalized dyadics where convenient.  Absolute value, cross-multiplied order,
addition, multiplication, and finite sums are primitive recursive on the
codes.  Denominator positivity is part of every rational-domain guard.

The relation \eqref{eq:regular-name} is a first-order schema for each fixed
evaluator symbol.  It does not quantify over functions.  Equality of two
names is expressed by
\[
 |a(p,k)-b(q,\ell)|\le 2^{-k}+2^{-\ell}
\]
for all \(k,\ell\), or by an equivalent one-precision formulation with a
constant shift.

\subsection{Addition, multiplication, and inverse}

Addition needs only a constant precision shift.  Define
\[
 c(p,q,k)=a(p,k+2)+b(q,k+2).
\]
Two applications of \eqref{eq:regular-name} and the triangle inequality show
that \(c\) is again regular.  Negation is immediate.

Multiplication additionally needs a magnitude contract.  If all queried
samples of \(a\) and \(b\) have absolute value at most \(M\), then
\[
 |a_kb_k-a_\ell b_\ell|
 \le M|a_k-a_\ell|+M|b_k-b_\ell|.
\]
A primitive-recursive shift \(r(M,k)\) making the right side at most the
requested dyadic error yields a multiplication name.  The magnitude bound may
be supplied globally or computed from a coarse sample plus the regularity
radius.

For inverse, suppose a positive rational \(d\) certifies
\(|a_k|\ge d\) at every precision queried by the algorithm.  Then
\[
 \left|\frac1{a_k}-\frac1{a_\ell}\right|
 =\frac{|a_k-a_\ell|}{|a_ka_\ell|}
 \le d^{-2}|a_k-a_\ell|.
\]
Choose the input precision so that the final expression meets the output
budget.  Numerically, inverse is totalized; correctness is asserted only under
the apartness guard.  Division is multiplication followed by guarded inverse.

\begin{theorem}[Arithmetic of guarded regular names]
\label{thm:regular-name-arithmetic}
For every fixed family of input evaluators, iPRA proves regularity and
extensionality of addition, negation, multiplication under a magnitude bound,
and inverse under a positive apartness bound.  The required precision
transformers are primitive recursive and compose.
\end{theorem}

\begin{proof}
Fix an output precision and instantiate the finitely many input regularity,
magnitude, and apartness facts used at that precision.  Clear positive
rational denominators and prove the resulting integer inequalities.  Each
calculation is a bounded trace whose invariant is the current rational error
budget.  Apply the trace compiler, then universally quantify the output
precision.  The displayed elementary inequalities give the local soundness
proofs.
\end{proof}

\subsection{Finite sums, products, and polynomial evaluation}

A finite sum carries the invariant
\[
 \left|\sum_{i<j}\widetilde x_i-
       \sum_{i<j}x_i\right|
 \le \sum_{i<j}e_i.
\]
A finite product carries a magnitude bound and the telescoping identity
\[
 \prod_{i<n}x_i-\prod_{i<n}y_i
 =\sum_{j<n}(x_j-y_j)
   \left(\prod_{i<j}x_i\right)
   \left(\prod_{j<i<n}y_i\right).
\]
Both are proved by quantifier-free induction on the finite length after all
error quantities have been rationally encoded.  Horner evaluation gives a
particularly compact polynomial trace and error recurrence.  Matrix
multiplication, determinants by a permutation sum or fraction-free
elimination, truncated power series, and finite convolutions are instances of
the same library.

\subsection{Finite-net transfer}

Let \(K(p)\) be a represented compact domain only in the following concrete
sense: for each rational radius \(\rho\), a primitive-recursive function
returns a finite list \(N(p,\rho)\), and a certificate proves
\[
 \forall x\in K(p)\;\exists y\in N(p,\rho)\; d(x,y)\le\rho.
\]
The quantifiers over represented points are interpreted through the finite
presentation used by the theorem; in a fully embedded statement, coverage is
a certificate-relative first-order formula.

Let \(f\) have a modulus \(\omega\): if \(d(x,y)\le\rho\), then
\(|f(x)-f(y)|\le\omega(\rho)\).

\begin{theorem}[Finite-net transfer]
\label{thm:finite-net-transfer}
If \(N\) is a \(\rho\)-net, \(f(y)\ge L\) for every \(y\in N\), and
\(\omega(\rho)\le\eta\), then \(f(x)\ge L-\eta\) throughout the represented
domain.  Likewise, a grid point whose value is within \(\eta\) of the finite
grid minimum is within \(2\eta\) of the global infimum.
\end{theorem}

\begin{proof}
For an arbitrary represented \(x\), the finite coverage certificate and
bounded search give a net point \(y\) with \(d(x,y)\le\rho\).  The modulus
gives \(f(x)\ge f(y)-\eta\ge L-\eta\).  For minimization, compare the selected
grid point to a net point near an arbitrary \(x\), then take the finite minimum
using bounded search.  No abstract compactness or completed infimum object is
used.
\end{proof}

The minimum \emph{value} needs no unique minimizer.  A coherent minimizer name
requires a quantitative uniqueness or shrinking-viability hypothesis.

\subsection{Riemann sums and integration}

An integral theorem is represented by a quadrature algorithm plus two finite
contracts:

\begin{enumerate}
\item a modulus controlling the change of the integrand within each cell; and
\item a tail or truncation estimate if the interval or series is unbounded.
\end{enumerate}

For a bounded interval and a mesh of width \(\rho\), upper and lower rational
sums differ by at most interval length times \(\omega(\rho)\).  The trace
computes the finite sum, while the invariant accumulates rounding and mesh
errors.  Additivity, linearity, and change of variables for explicitly
monotone rational maps reduce to finite sum rearrangements and modulus bounds.

The fundamental theorem of calculus should therefore be split into finite
statements.  For the derivative-to-integral direction, telescope finite
increments and control each local remainder.  For the integral-to-derivative
direction, compare a difference quotient to cell averages using a continuity
modulus.  A theorem packet that states merely that a derivative is continuous
has not yet supplied the numerical input required by the iPRA proof.

\subsection{Series and analytic functions}

A series-defined function is a fixed finite-sum evaluator together with a
primitive-recursive tail bound
\[
 \left|\sum_{j\ge N}t_j(p,z)\right|\le T(p,z,N).
\]
All algebraic identities are first proved for finite truncations.  Tail bounds
then transfer them to the represented function at a fixed requested
precision.  This pattern covers exponential, sine, cosine, logarithm on a
guarded domain, Euler products over bounded primes, and many special-function
identities.

Local analytic inversion uses a finite certificate: a center, rational box,
derivative or Jacobian gap, contraction constant, coefficient or derivative
bounds, and a radius.  Finite iteration constructs approximations; the
contraction estimate gives horizontal shrinking.  Singular branch discovery
is a separate producer problem and must not be hidden in the local theorem.

\subsection{Approximate extrema}

For a finite rational grid, bounded search returns a least residual or least
value point.  Let \(m_k\) denote the finite minimum.  The finite-net theorem
turns \(m_k\) into a name for the infimum value.  To obtain a point:

\begin{itemize}
\item retain the relation of all \(\eta_k\)-near-minimizers; or
\item provide a quantitative uniqueness estimate, such as strong convexity,
      and apply shrinking viability.
\end{itemize}

This resolves a common constructive mismatch.  The operation
``take the infimum value'' is often stable, while ``take the unique point where
it is attained'' needs more data.

\subsection{Finite probability laws}

A finite rational law is a list of outcomes and nonnegative rational weights
summing to one.  Expectation, variance, characteristic functions at rational
arguments, convolution, and distribution values at rational thresholds are
finite sums.  Approximation of an unbounded or continuous law requires
explicit tail and quantization moduli.  Approximation of a discontinuous CDF
at a threshold additionally requires an anti-concentration or boundary-mass
bound.

Thus a constructive central-limit statement has a transparent shape:
finite rational laws, a finite convolution or characteristic-function trace,
a quantitative error estimate, and a rational threshold protected by a
boundary budget.  Higher-type measures are unnecessary unless they are given
through a fixed finite-observation evaluator with the corresponding moduli.

\subsection{Precision ledgers}

Every analytic module should export a machine-readable precision ledger:

\begin{longtable}{@{}p{0.24\textwidth}p{0.68\textwidth}@{}}
\toprule
Field & Meaning\\
\midrule
\endhead
Output challenge & Requested error, region, derivative order, or test.\\
Input queries & Exact finite list or adaptive trace of samples.\\
Magnitude data & Bounds needed for products, tails, or matrix operations.\\
Apartness data & Positive margins needed by inverses and exact branches.\\
Truncation bound & Number of series terms, grid cells, or iterations.\\
Roundoff budget & Rational allocation to every finite operation.\\
Analytic remainder & Tail, modulus, contraction, or residual estimate.\\
Final inequality & Arithmetic implication proving the output contract.\\
\bottomrule
\end{longtable}

The ledger is not auxiliary documentation.  It is the numerical content from
which the fixed-precision object-language proof is built.
